% chapters/thyfermi.tex
% generated from tiristore_bande_fermi.tex
\chapter{Fermi levels in a thyristor}
\label{ch:thyfermi}
\vspace{-12mm}

\begin{chapabstract}
\noindent
Il documento descrive il funzionamento del tiristore (SCR) dal punto di vista della
fisica del semiconduttore, usando come strumento principale il diagramma a bande di
energia con il livello di Fermi e i quasi-livelli di Fermi. Vengono trattati
l'equilibrio termodinamico, il blocco diretto, l'innesco e lo stato di conduzione.
Le figure sono qualitative e non in scala: le tensioni reali di blocco (kV)
non sono rappresentabili sulla stessa scala del gap del silicio ($E_g \simeq 1{,}12$~eV).
\end{chapabstract}

% =====================================================================
\section{Struttura del dispositivo}
% =====================================================================

Il tiristore è una struttura a quattro strati $p_1$--$n_1$--$p_2$--$n_2$ con tre
giunzioni $J_1$, $J_2$, $J_3$ e tre terminali: anodo (A) sull'emettitore $p_1$,
catodo (K) sull'emettitore $n_2$, gate (G) sulla base $p_2$.

\begin{figure}[h]
\centering
\begin{tikzpicture}[x=1cm,y=1cm]
  \fill[pcol] (0,0) rectangle (2.5,2);
  \fill[ncol] (2.5,0) rectangle (7,2);
  \fill[pcol] (7,0) rectangle (9.5,2);
  \fill[ncol] (9.5,0) rectangle (12,2);
  \draw[thick] (0,0) rectangle (12,2);
  \draw[thick] (2.5,0)--(2.5,2) (7,0)--(7,2) (9.5,0)--(9.5,2);
  \node at (1.25,1.35) {$p_1$};  \node[font=\scriptsize] at (1.25,0.8) {$\sim 10^{19}\,\mathrm{cm^{-3}}$};
  \node at (4.75,1.35) {$n_1$};  \node[font=\scriptsize] at (4.75,0.8) {$\sim 10^{13}\,\mathrm{cm^{-3}}$ (base larga)};
  \node at (8.25,1.35) {$p_2$};  \node[font=\scriptsize] at (8.25,0.8) {$\sim 10^{17}\,\mathrm{cm^{-3}}$};
  \node at (10.75,1.35) {$n_2$}; \node[font=\scriptsize] at (10.75,0.8) {$\sim 10^{19}\,\mathrm{cm^{-3}}$};
  \node[font=\small] at (2.5,2.3) {$J_1$};
  \node[font=\small] at (7,2.3)   {$J_2$};
  \node[font=\small] at (9.5,2.3) {$J_3$};
  \draw[very thick] (0,1)--(-1,1) node[left] {A};
  \draw[very thick] (12,1)--(13,1) node[right] {K};
  \draw[very thick] (8.25,0)--(8.25,-1) node[below] {G};
\end{tikzpicture}
\caption{Struttura $p$--$n$--$p$--$n$ e livelli di drogaggio tipici di un tiristore
di potenza per controllo di fase.}
\end{figure}

Il ruolo dei singoli strati:
\begin{itemize}\itemsep2pt
  \item $p_1$ ed $n_2$ sono \emph{emettitori} fortemente drogati: iniettano
        rispettivamente lacune ed elettroni con alta efficienza;
  \item $n_1$ è la \emph{base larga} debolmente drogata: è lo strato che sostiene
        la tensione di blocco e che, in conduzione, viene invaso dal plasma
        elettroni--lacune;
  \item $p_2$ è la \emph{base di controllo}, collegata al gate.
\end{itemize}

% =====================================================================
\section{Livello di Fermi e quasi-livelli di Fermi}
% =====================================================================

All'equilibrio termodinamico le concentrazioni di portatori sono legate a un
\emph{unico} livello di Fermi $E_F$:
\begin{equation}
n = n_i \exp\!\left(\frac{E_F-E_i}{kT}\right),\qquad
p = n_i \exp\!\left(\frac{E_i-E_F}{kT}\right),\qquad
np = n_i^2 .
\end{equation}
La condizione di equilibrio è $\nabla E_F = 0$: \textbf{il livello di Fermi è una
retta orizzontale su tutto il dispositivo}. Le bande $E_C$ ed $E_V$ si piegano in
corrispondenza delle giunzioni proprio perché $E_F$ deve restare piatto: la
piegatura corrisponde al potenziale di built-in
\begin{equation}
V_{bi}=\frac{kT}{q}\,\ln\!\left(\frac{N_A N_D}{n_i^2}\right).
\end{equation}

Fuori equilibrio elettroni e lacune non sono più in equilibrio tra loro e servono
due potenziali elettrochimici distinti, i \emph{quasi-livelli di Fermi} $E_{Fn}$
ed $E_{Fp}$ (imref):
\begin{equation}
n = n_i \exp\!\left(\frac{E_{Fn}-E_i}{kT}\right),\qquad
p = n_i \exp\!\left(\frac{E_i-E_{Fp}}{kT}\right),\qquad
np = n_i^2 \exp\!\left(\frac{E_{Fn}-E_{Fp}}{kT}\right).
\end{equation}
Due relazioni sono essenziali per leggere i diagrammi che seguono:
\begin{equation}
\boxed{\;J_n = \mu_n\,n\,\nabla E_{Fn}\;,\qquad J_p = \mu_p\,p\,\nabla E_{Fp}\;}
\label{thyfermi:eq:corrente}
\end{equation}
cioè \emph{la corrente scorre solo dove il quasi-livello di Fermi ha pendenza};
e la separazione dei quasi-livelli dice se localmente c'è iniezione
($E_{Fn}>E_{Fp}$, $np>n_i^2$) o svuotamento ($E_{Fn}<E_{Fp}$, $np<n_i^2$).
Infine, la tensione applicata ai morsetti è la differenza dei livelli di Fermi
dei due contatti:
\begin{equation}
qV_{AK} = E_F(\mathrm{K}) - E_F(\mathrm{A}).
\end{equation}

% =====================================================================
\section{Tiristore a riposo: equilibrio termodinamico ($V_{AK}=0$)}
% =====================================================================

Con i morsetti in cortocircuito non circola corrente: $E_F$ è unico e piatto.
Le tre giunzioni sono tutte all'equilibrio e presentano ciascuna la propria
barriera di built-in. Il profilo delle bande è quello di figura~\ref{thyfermi:fig:eq}:
scendono attraversando $J_1$ ($p_1 \to n_1$), risalgono attraversando $J_2$
($n_1 \to p_2$) e riscendono attraversando $J_3$ ($p_2 \to n_2$).

\begin{figure}[h]
\centering
\begin{tikzpicture}[x=1cm,y=2cm]
  \bandbg{-1.35}{1.25}
  % banda di conduzione
  \draw[ccol,very thick]
    (0,1.07) -- (2.15,1.07) .. controls (2.5,1.07) and (2.5,0.30) .. (2.85,0.30)
    -- (6.60,0.30) .. controls (7,0.30) and (7,0.92) .. (7.40,0.92)
    -- (9.15,0.92) .. controls (9.5,0.92) and (9.5,0.05) .. (9.85,0.05)
    -- (12,0.05);
  % banda di valenza
  \draw[vcol,very thick]
    (0,-0.05) -- (2.15,-0.05) .. controls (2.5,-0.05) and (2.5,-0.82) .. (2.85,-0.82)
    -- (6.60,-0.82) .. controls (7,-0.82) and (7,-0.20) .. (7.40,-0.20)
    -- (9.15,-0.20) .. controls (9.5,-0.20) and (9.5,-1.07) .. (9.85,-1.07)
    -- (12,-1.07);
  % livello di Fermi
  \draw[black,thick,dashed] (0,0) -- (12,0);
  \node[anchor=south,font=\small,ccol] at (3.6,0.30) {$E_C$};
  \node[anchor=north,font=\small]      at (3.6,0.00) {$E_F$};
  \node[anchor=north,font=\small,vcol] at (3.6,-0.82) {$E_V$};
  \draw[<->,>=Stealth] (5.6,0.30) -- (5.6,-0.82)
        node[midway,fill=white,inner sep=1pt,font=\scriptsize] {$E_g$};
  \draw[<->,>=Stealth,gray] (7.05,0.30) -- (7.05,0.92)
        node[midway,right=2pt,fill=white,inner sep=1pt,font=\scriptsize] {$qV_{bi,J_2}$};
\end{tikzpicture}
\caption{Diagramma a bande all'equilibrio ($V_{AK}=0$). Il livello di Fermi è unico
e orizzontale; le bande si piegano alle tre giunzioni. Figura qualitativa, non in scala.}
\label{thyfermi:fig:eq}
\end{figure}

\noindent
Da \eqref{thyfermi:eq:corrente}: $E_F$ piatto $\Rightarrow$ $J_n=J_p=0$. In ogni giunzione
la corrente di diffusione è esattamente compensata da quella di deriva.

% =====================================================================
\section{Stato di blocco diretto ($V_{AK}>0$, gate aperto)}
% =====================================================================

Applicando l'anodo positivo, $J_1$ e $J_3$ risultano polarizzate direttamente
mentre $J_2$ è polarizzata inversamente. Poiché la corrente che può circolare
è limitata dalla giunzione inversa, quasi tutta la tensione applicata cade su
$J_2$: le due giunzioni dirette restano a poche decine di millivolt.

La regione di carica spaziale di $J_2$ si estende quasi esclusivamente nella base
$n_1$, che è debolmente drogata: è questo il motivo per cui $n_1$ deve essere
larga (centinaia di $\mu$m) in un dispositivo da kV.

\begin{figure}[h]
\centering
\begin{tikzpicture}[x=1cm,y=1.55cm]
  \bandbg{-2.45}{1.15}
  \fill[black,opacity=0.06] (5.4,-2.45) rectangle (7.35,1.15);
  % banda di conduzione
  \draw[ccol,very thick]
    (0,-0.33) -- (2.15,-0.33) .. controls (2.5,-0.33) and (2.5,-1.10) .. (2.85,-1.10)
    -- (5.40,-1.10) .. controls (6.40,-1.10) and (6.90,0.92) .. (7.35,0.92)
    -- (9.15,0.92) .. controls (9.5,0.92) and (9.5,0.05) .. (9.85,0.05)
    -- (12,0.05);
  % banda di valenza
  \draw[vcol,very thick]
    (0,-1.45) -- (2.15,-1.45) .. controls (2.5,-1.45) and (2.5,-2.22) .. (2.85,-2.22)
    -- (5.40,-2.22) .. controls (6.40,-2.22) and (6.90,-0.20) .. (7.35,-0.20)
    -- (9.15,-0.20) .. controls (9.5,-0.20) and (9.5,-1.07) .. (9.85,-1.07)
    -- (12,-1.07);
  % quasi-livelli di Fermi
  \draw[fncol,thick,dashed]
    (0,-1.40) -- (7.30,-1.40) .. controls (7.80,-1.40) and (7.90,0) .. (8.40,0) -- (12,0);
  \draw[fpcol,thick,dash dot]
    (12,0) -- (5.60,0) .. controls (5.10,0) and (5.00,-1.40) .. (4.50,-1.40) -- (0,-1.40);
  \node[anchor=south,font=\small,ccol] at (3.6,-1.10) {$E_C$};
  \node[anchor=north,font=\small,vcol] at (3.6,-2.22) {$E_V$};
  \node[anchor=south,font=\small,fncol] at (9.9,0.02) {$E_{Fn}$};
  \node[anchor=north,font=\small,fpcol] at (2.6,-1.42) {$E_{Fp}$};
  \draw[<->,>=Stealth] (6.0,-1.40) -- (6.0,0)
        node[pos=0.88,fill=white,inner sep=1pt,font=\scriptsize] {$qV_{AK}$};
  \node[font=\scriptsize,align=center] at (6.35,-2.15) {zona di\\svuotamento};
\end{tikzpicture}
\caption{Blocco diretto: $J_2$ è polarizzata inversamente e sostiene tutta la tensione.
I quasi-livelli di Fermi si separano di $qV_{AK}$ attraverso la zona svuotata, con
$E_{Fn}<E_{Fp}$ ($np<n_i^2$: generazione netta).}
\label{thyfermi:fig:blocco}
\end{figure}

\noindent
Osservazioni sulla figura~\ref{thyfermi:fig:blocco}:
\begin{itemize}\itemsep2pt
  \item $E_{Fn}$ è agganciato al contatto d'anodo, $E_{Fp}$ a quello di catodo;
        la loro separazione nella zona svuotata vale esattamente $qV_{AK}$;
  \item essendo $E_{Fn}<E_{Fp}$ si ha $np<n_i^2$: nella zona svuotata prevale la
        \emph{generazione} termica (Shockley--Read--Hall), che produce la corrente
        di fuga $I_{D}$ del dispositivo, fortemente dipendente dalla temperatura;
  \item entrambi i quasi-livelli sono quasi piatti nelle regioni neutre: lì la
        corrente è piccola, coerentemente con \eqref{thyfermi:eq:corrente}.
\end{itemize}

\subsection{Modello a due transistor}

La struttura a quattro strati si interpreta come due transistor bipolari
accoppiati: $Q_1$ ($p_1n_1p_2$, pnp, guadagno $\alpha_1$) e $Q_2$ ($n_1p_2n_2$,
npn, guadagno $\alpha_2$), ciascuno che alimenta la base dell'altro.

\begin{figure}[h]
\centering
\begin{circuitikz}[scale=1.0,transform shape]
  \draw (0,0.8)  node[pnp,xscale=-1] (Q1) {};
  \draw (2.2,-0.8) node[npn] (Q2) {};
  \node[font=\small] at (-0.9,0.8) {$Q_1$ (pnp), $\alpha_1$};
  \node[font=\small] at (3.4,-0.8) {$Q_2$ (npn), $\alpha_2$};
  \draw (Q1.E) -- (0,2.0) node[above] {A};
  \draw (Q2.E) -- (2.2,-2.4) node[below] {K};
  \draw (Q1.C) -- (0,-0.8) -- (Q2.B);
  \draw (Q2.C) -- (2.2,0.8) -- (Q1.B);
  \draw (0.85,-0.8) node[circle,fill,inner sep=1pt]{} -- (0.85,-2.0) node[below] {G};
\end{circuitikz}
\caption{Modello a due transistor accoppiati. La base $n_1$ di $Q_1$ è il
collettore di $Q_2$ e viceversa.}
\end{figure}

Imponendo la continuità della corrente si ottiene
\begin{equation}
I_A=\frac{\alpha_2 I_G + I_{CO1}+I_{CO2}}{1-(\alpha_1+\alpha_2)} .
\label{thyfermi:eq:duetr}
\end{equation}
Finché $\alpha_1+\alpha_2 \ll 1$ il dispositivo blocca. Nei tiristori i guadagni
$\alpha$ crescono con la corrente (per bassa iniezione l'efficienza di emettitore
e il fattore di trasporto sono ridotti), quindi la condizione
$\alpha_1+\alpha_2 \to 1$ è raggiungibile aumentando la corrente di gate.

% =====================================================================
\section{Innesco}
% =====================================================================

La sequenza fisica dell'innesco, letta sulle bande:
\begin{enumerate}\itemsep2pt
  \item la corrente di gate inietta lacune in $p_2$, che polarizza direttamente
        $J_3$;
  \item l'emettitore $n_2$ inietta elettroni in $p_2$; questi diffondono e vengono
        \emph{raccolti} dal campo di $J_2$, che li accelera nella base $n_1$;
  \item gli elettroni accumulati in $n_1$ abbassano il potenziale della base e
        polarizzano direttamente $J_1$: l'emettitore $p_1$ inietta lacune;
  \item le lacune attraversano $n_1$, raggiungono $p_2$ e sostengono
        ulteriormente l'iniezione di $J_3$: la retroazione è positiva.
\end{enumerate}
Quando $\alpha_1+\alpha_2=1$ il denominatore di \eqref{thyfermi:eq:duetr} si annulla: la
corrente non è più determinata dal dispositivo ma dal circuito esterno, e il gate
perde il controllo. In termini di diagramma a bande, la barriera di $J_2$ si
abbassa progressivamente man mano che la base $n_1$ si riempie di plasma, finché
il massimo di $E_C$ in corrispondenza di $J_2$ scompare del tutto.

Perché lo stato di conduzione si mantenga occorre che la corrente superi la
corrente di aggancio $I_L$ (latching) durante l'innesco e resti sopra la corrente
di mantenimento $I_H$ (holding) durante la conduzione.

% =====================================================================
\section{Stato di conduzione}
% =====================================================================

In conduzione \textbf{tutte e tre le giunzioni sono polarizzate direttamente}.
La base $n_1$ è invasa da un plasma quasi neutro con $n \simeq p \gg N_D$
(alta iniezione): la sua resistività crolla di ordini di grandezza rispetto al
valore di equilibrio. È la \emph{modulazione della conducibilità}, che consente a
uno strato spesso e poco drogato di condurre migliaia di ampere con una caduta di
1--2~V.

\begin{figure}[h]
\centering
\begin{tikzpicture}[x=1cm,y=2.5cm]
  \bandbg{-1.30}{0.72}
  % banda di conduzione
  \draw[ccol,very thick]
    (0,0.45) -- (2.15,0.45) .. controls (2.5,0.45) and (2.5,0.25) .. (2.85,0.25)
    -- (6.60,0.25) .. controls (7,0.25) and (7,0.35) .. (7.40,0.35)
    -- (9.15,0.35) .. controls (9.5,0.35) and (9.5,0.05) .. (9.85,0.05)
    -- (12,0.05);
  % banda di valenza
  \draw[vcol,very thick]
    (0,-0.67) -- (2.15,-0.67) .. controls (2.5,-0.67) and (2.5,-0.87) .. (2.85,-0.87)
    -- (6.60,-0.87) .. controls (7,-0.87) and (7,-0.77) .. (7.40,-0.77)
    -- (9.15,-0.77) .. controls (9.5,-0.77) and (9.5,-1.07) .. (9.85,-1.07)
    -- (12,-1.07);
  % quasi-livelli di Fermi
  \draw[fncol,thick,dashed]
    (12,0) -- (2.90,0) .. controls (2.40,0) and (2.30,-0.62) .. (1.80,-0.62) -- (0,-0.62);
  \draw[fpcol,thick,dash dot]
    (0,-0.62) -- (9.90,-0.62) .. controls (10.40,-0.62) and (10.50,0) .. (11.00,0) -- (12,0);
  % plasma
  \foreach \x in {3.4,4.1,4.8,5.5,6.2} {
    \fill[ccol] (\x,0.16) circle (1.1pt);
    \draw[vcol] (\x+0.25,-0.78) circle (1.1pt);
  }
  \node[anchor=south,font=\small,ccol] at (3.3,0.25) {$E_C$};
  \node[anchor=north,font=\small,vcol] at (3.3,-0.87) {$E_V$};
  \node[anchor=south,font=\small,fncol] at (8.3,0.01) {$E_{Fn}$};
  \node[anchor=north,font=\small,fpcol] at (8.3,-0.63) {$E_{Fp}$};
  \draw[<->,>=Stealth] (7.6,0) -- (7.6,-0.62)
        node[midway,fill=white,inner sep=1pt,font=\scriptsize] {$E_{Fn}\!-\!E_{Fp}$};
\end{tikzpicture}
\caption{Stato di conduzione: le barriere sono quasi annullate, $E_{Fn}$ è agganciato
al catodo ed $E_{Fp}$ all'anodo. Nella base $n_1$ i due quasi-livelli sono separati e
vicini alle rispettive bande: $np \gg n_i^2$, plasma di elettroni ($\bullet$) e
lacune ($\circ$).}
\label{thyfermi:fig:on}
\end{figure}

\noindent
Come leggere la figura~\ref{thyfermi:fig:on}:
\begin{itemize}\itemsep2pt
  \item $E_{Fn}$ resta praticamente piatto dal catodo fino a $p_1$ e $E_{Fp}$
        dall'anodo fino a $n_2$: ciascun tipo di portatore attraversa quasi tutto
        il dispositivo come portatore \emph{minoritario} in eccesso;
  \item nella base la separazione $E_{Fn}-E_{Fp}$ è grande e le distanze
        $E_C-E_{Fn}$ ed $E_{Fp}-E_V$ sono piccole e confrontabili: è la firma
        dell'alta iniezione ambipolare;
  \item la lieve pendenza dei quasi-livelli nella base rappresenta la caduta
        ohmica residua sul plasma;
  \item il massimo di $E_C$ a $J_2$ è scomparso: la giunzione di blocco non è più
        tale.
\end{itemize}

La tensione ai morsetti è la somma algebrica delle tre cadute di giunzione più il
termine resistivo. Poiché $J_2$ è orientata in senso opposto a $J_1$ e $J_3$:
\begin{equation}
V_T = V_{J_1} - V_{J_2} + V_{J_3} + R_{\text{mod}} I_A \;\simeq\; 1\div 2~\text{V}.
\end{equation}
Questo spiega perché la caduta di un tiristore in conduzione è maggiore di quella
di un diodo ma molto minore della somma di tre giunzioni.

% =====================================================================
\section{Blocco inverso e spegnimento}
% =====================================================================

Con $V_{AK}<0$ le giunzioni $J_1$ e $J_3$ sono polarizzate inversamente. $J_3$ ha
entrambi i lati fortemente drogati e va in valanga a poche decine di volt: è
quindi $J_1$, affacciata sulla base $n_1$, a sostenere la tensione inversa. Il
diagramma è speculare a quello di figura~\ref{thyfermi:fig:blocco}, con la zona svuotata
che si estende da $J_1$ verso destra e $E_{Fp}$ agganciato all'anodo, più basso
di $E_{Fn}$.

Lo spegnimento richiede due fasi distinte, entrambe visibili nel quadro a bande:
\begin{enumerate}\itemsep2pt
  \item \textbf{rimozione del plasma}. La carica immagazzinata nella base $n_1$
        deve essere estratta dalla corrente inversa e/o ricombinare con costante
        di tempo pari alla vita media $\tau$ dei portatori. È l'origine della
        carica di recupero inverso $Q_{rr}$ e del picco $I_{RM}$.
  \item \textbf{recupero della barriera di $J_2$}. Solo quando la concentrazione
        residua nella base è scesa abbastanza i quasi-livelli tornano a
        riseparare $E_{Fn}<E_{Fp}$ attorno a $J_2$ e la giunzione ricostruisce la
        propria zona svuotata. Il tempo necessario è il \emph{tempo di spegnimento
        del circuito} $t_q$: se si riapplica tensione diretta prima di $t_q$, la
        carica residua riporta $\alpha_1+\alpha_2$ a 1 e il dispositivo si
        riaccende senza comando di gate.
\end{enumerate}
Questa è la ragione fisica per cui un modello circuitale accurato del tiristore
non può limitarsi alla caratteristica statica, ma deve includere una dinamica di
carica (per esempio del tipo a carica concentrata di Lauritzen--Ma), con la vita
media $\tau$ e il tempo di transito come parametri.

% =====================================================================
\section{Riepilogo}
% =====================================================================

\begin{table}[h]
\centering
\small
\begin{tabular}{@{}lccccl@{}}
\toprule
Stato & $J_1$ & $J_2$ & $J_3$ & Livelli di Fermi & Caduta \\
\midrule
Equilibrio        & --      & --      & --      & $E_F$ unico e piatto & $0$ \\
Blocco diretto    & diretta & inversa & diretta & $E_{Fn}<E_{Fp}$ su $J_2$ & tutta su $J_2$ \\
Blocco inverso    & inversa & diretta & inversa & $E_{Fp}<E_{Fn}$ su $J_1$ & tutta su $J_1$ \\
Conduzione        & diretta & diretta & diretta & $E_{Fn}>E_{Fp}$ nella base & $1\div2$~V \\
\bottomrule
\end{tabular}
\caption{Sintesi degli stati operativi.}
\end{table}

\paragraph{Regola di lettura.} Un solo livello di Fermi piatto significa equilibrio
e corrente nulla. Due quasi-livelli separati con $E_{Fn}<E_{Fp}$ significano
svuotamento e generazione: il dispositivo blocca. Due quasi-livelli separati con
$E_{Fn}>E_{Fp}$ significano iniezione e ricombinazione: il dispositivo conduce.

% =====================================================================
\section*{Riferimenti essenziali}
\begin{itemize}\itemsep1pt
  \item B.\,J. Baliga, \emph{Fundamentals of Power Semiconductor Devices},
        Springer, cap.\ sui tiristori.
  \item S.\,M. Sze, K.\,K. Ng, \emph{Physics of Semiconductor Devices},
        Wiley, cap.\ ``Thyristors''.
  \item J. Lutz et al., \emph{Semiconductor Power Devices}, Springer.
  \item P.\,O. Lauritzen, C.\,L. Ma, ``A simple diode model with reverse recovery'',
        \emph{IEEE Trans.\ Power Electron.}, 1991.
\end{itemize}

